% Source: https://artofproblemsolving.com/community/c3h2290078_modular_arithmetic_troll
% Stable identity: aops-topic-2290078-q02
Since $196=4\cdot49$ with $\gcd(4,49)=1$, we need $4\mid x(x-3)$ and $49\mid x(x-3)$.

\textbf{Modulo $4$.} Checking $x=0,1,2,3$: $x(x-3)\equiv0,2,2,0\pmod4$. So $x\equiv0$ or $3\pmod4$.

\textbf{Modulo $49$.} A common divisor of $x$ and $x-3$ divides $3$, so $7$ cannot divide both. Hence if $49\mid x(x-3)$, both factors of $7$ come from the same factor: $49\mid x$ or $49\mid x-3$. So $x\equiv0$ or $3\pmod{49}$.

\textbf{Combining} (Chinese remainder theorem), there are four residues modulo $196$:
\begin{itemize}
\item $x\equiv0\pmod4$, $x\equiv0\pmod{49}$: $x\equiv0$;
\item $x\equiv3\pmod4$, $x\equiv3\pmod{49}$: $x\equiv3$;
\item $x\equiv0\pmod4$, $x\equiv3\pmod{49}$: $x=3+49k$ with $3+k\equiv0\pmod4$, so $k=1$ and $x\equiv52$;
\item $x\equiv3\pmod4$, $x\equiv0\pmod{49}$: $x=49k$ with $k\equiv3\pmod4$, so $x\equiv147$.
\end{itemize}
The answer is $\boxed{x\equiv0,\ 3,\ 52,\ 147\pmod{196}}$.
